\documentclass[11pt]{article}
\usepackage[margin=2.2cm]{geometry}
\usepackage{graphicx}
\usepackage{booktabs}
\usepackage{amsmath,amssymb}
\usepackage{hyperref}
\usepackage{xcolor}
\usepackage{caption}
\usepackage{float}
\newcommand{\Rsq}{R^2}

\title{PySR Warm-Start (Hybrid) vs.\ Cold PySR:\\
Trajectory Analysis Report Across 5 Seeds}
\author{Generated by \texttt{generate\_figures.py} (exp3 trajectory block)}
\date{\today}

\begin{document}
\maketitle

\begin{abstract}
\noindent\textbf{Important data-provenance caveat first.} This report was generated from a
\emph{synthetic} 5-seed $\times$ 12-equation demonstration dataset built to validate the
trajectory-analysis pipeline (\texttt{generate\_figures.py}, exp3 block), \emph{not} from real
HypatiaX/PySR experiment logs. No real \texttt{exp3\_nguyen12\_seed*.json} result files were
available in this conversation. Every number and figure below is real \emph{output of the
pipeline} running on this synthetic input, and the pipeline logic itself is unchanged from what
would run on real data --- but the specific conclusions (win rates, speedup factors, $p$-values)
describe the synthetic data's built-in bias, not the actual HypatiaX/PySR-only comparison. They
must be regenerated by pointing \texttt{--results-dir} at the real seed files before being cited
anywhere, including the paper.

\smallskip
\noindent With that caveat stated, the \emph{methodological} conclusion is real and does not
depend on which dataset is used: two plausible metrics for ``does the hybrid warm-start help''
can disagree, and this run reproduces that disagreement by construction. Using \textbf{final loss
at whatever iteration each arm's run happened to stop}, the cold-PySR arm (P) appears better in
52/60 equation--seed pairs (Wilcoxon $p = 2.0\times10^{-8}$) --- but this arm was simply allowed to
run longer, so this is not evidence about the warm-start itself. Using \textbf{iterations needed
to first reach the paper's own success threshold} ($\Rsq \ge 0.9999$), the hybrid arm (H) is
faster in 60/60 equation--seed pairs (Wilcoxon $p = 1.6\times10^{-11}$), a pooled geometric-mean
speedup of $3.38\times$. We recommend the threshold-crossing speedup as the primary metric for any
``warm-start works'' claim, and flag final-loss comparisons as requiring a matched iteration/time
budget before they mean anything. This same discipline --- computing summary statistics
programmatically from source data rather than by hand, and stating exactly what was measured ---
is also the fix needed for the unrelated Nguyen-12 success-count arithmetic issue flagged
separately in \texttt{01\_nguyen12\_success\_count\_arithmetic.md} (three disagreeing
success-count denominators for the same table), and the same regenerate-from-source approach
should be applied there.
\end{abstract}

\tableofcontents

\section{Methodology}
\label{sec:methodology}

Each seed file records, per Nguyen-12 equation and per arm (H = LLM warm-start hybrid, P = cold
PySR), a polled trajectory of $(\text{iteration}, \text{best\_loss}, \text{elapsed\_seconds})$
observations, plus a \texttt{trajectory\_summary} with \texttt{final\_best\_loss} and
\texttt{first\_threshold\_iteration} (first iteration at which $\Rsq \ge 0.9999$ was observed).
For each seed we produce:

\begin{itemize}
  \item \textbf{Overview} --- the two arms' loss trajectories aggregated across all 12 equations
    into a single median curve per arm (H, P), rank-aligned by observation order with a 25th--75th
    percentile band, on a shared log-loss axis.
  \item \textbf{Per-equation grid} --- a $3\times4$ grid of subplots, one per equation, each on its
    own y-axis, showing the real (unaggregated) H and P loss trajectories, annotated with which arm
    reached the lower final loss and which reached the $\Rsq\ge0.9999$ threshold first.
  \item \textbf{Win/loss scatter} --- P's final loss (x) vs.\ H's final loss (y), log--log, per
    equation; points below the $y=x$ line are equations where H ended with a lower loss.
  \item \textbf{Speedup bar} --- for equations where both arms reached the threshold, the ratio
    (P's threshold iteration) / (H's threshold iteration); bars above $1\times$ favor H.
\end{itemize}

\noindent After processing all seed files, the same two paired comparisons (final loss, and
threshold-crossing speedup) are pooled across every seed and equation and tested with a paired
Wilcoxon signed-rank test (\texttt{scipy.stats.wilcoxon}), producing the two cross-seed figures in
Section~\ref{sec:allseeds}.

\section{Per-Seed Results}
\label{sec:perseed}

\FloatBarrier
%=====================================================================
\subsection{Seed 1}
\begin{figure}[H]\centering
\includegraphics[width=0.85\linewidth]{figures/exp3_traj_seed1_overview.png}
\caption{Seed 1 --- overview (median H vs.\ P loss trajectory, rank-aligned, IQR band).}
\end{figure}
\begin{figure}[H]\centering
\includegraphics[width=\linewidth]{figures/exp3_traj_seed1_grid.png}
\caption{Seed 1 --- per-equation trajectory grid (real, unaggregated H/P curves).}
\end{figure}
\begin{figure}[H]\centering
\includegraphics[width=0.68\linewidth]{figures/exp3_traj_seed1_winloss_scatter.png}
\caption{Seed 1 --- final-loss win/loss scatter.}
\end{figure}
\begin{figure}[H]\centering
\includegraphics[width=0.85\linewidth]{figures/exp3_traj_seed1_speedup_bar.png}
\caption{Seed 1 --- per-equation speedup to $\Rsq\ge0.9999$.}
\end{figure}

\FloatBarrier
%=====================================================================
\subsection{Seed 2}
\begin{figure}[H]\centering
\includegraphics[width=0.85\linewidth]{figures/exp3_traj_seed2_overview.png}
\caption{Seed 2 --- overview (median H vs.\ P loss trajectory, rank-aligned, IQR band).}
\end{figure}
\begin{figure}[H]\centering
\includegraphics[width=\linewidth]{figures/exp3_traj_seed2_grid.png}
\caption{Seed 2 --- per-equation trajectory grid (real, unaggregated H/P curves).}
\end{figure}
\begin{figure}[H]\centering
\includegraphics[width=0.68\linewidth]{figures/exp3_traj_seed2_winloss_scatter.png}
\caption{Seed 2 --- final-loss win/loss scatter.}
\end{figure}
\begin{figure}[H]\centering
\includegraphics[width=0.85\linewidth]{figures/exp3_traj_seed2_speedup_bar.png}
\caption{Seed 2 --- per-equation speedup to $\Rsq\ge0.9999$.}
\end{figure}

\FloatBarrier
%=====================================================================
\subsection{Seed 3}
\begin{figure}[H]\centering
\includegraphics[width=0.85\linewidth]{figures/exp3_traj_seed3_overview.png}
\caption{Seed 3 --- overview (median H vs.\ P loss trajectory, rank-aligned, IQR band).}
\end{figure}
\begin{figure}[H]\centering
\includegraphics[width=\linewidth]{figures/exp3_traj_seed3_grid.png}
\caption{Seed 3 --- per-equation trajectory grid (real, unaggregated H/P curves).}
\end{figure}
\begin{figure}[H]\centering
\includegraphics[width=0.68\linewidth]{figures/exp3_traj_seed3_winloss_scatter.png}
\caption{Seed 3 --- final-loss win/loss scatter.}
\end{figure}
\begin{figure}[H]\centering
\includegraphics[width=0.85\linewidth]{figures/exp3_traj_seed3_speedup_bar.png}
\caption{Seed 3 --- per-equation speedup to $\Rsq\ge0.9999$.}
\end{figure}

\FloatBarrier
%=====================================================================
\subsection{Seed 7}
\begin{figure}[H]\centering
\includegraphics[width=0.85\linewidth]{figures/exp3_traj_seed7_overview.png}
\caption{Seed 7 --- overview (median H vs.\ P loss trajectory, rank-aligned, IQR band).}
\end{figure}
\begin{figure}[H]\centering
\includegraphics[width=\linewidth]{figures/exp3_traj_seed7_grid.png}
\caption{Seed 7 --- per-equation trajectory grid (real, unaggregated H/P curves).}
\end{figure}
\begin{figure}[H]\centering
\includegraphics[width=0.68\linewidth]{figures/exp3_traj_seed7_winloss_scatter.png}
\caption{Seed 7 --- final-loss win/loss scatter.}
\end{figure}
\begin{figure}[H]\centering
\includegraphics[width=0.85\linewidth]{figures/exp3_traj_seed7_speedup_bar.png}
\caption{Seed 7 --- per-equation speedup to $\Rsq\ge0.9999$.}
\end{figure}

\FloatBarrier
%=====================================================================
\subsection{Seed 42}
\begin{figure}[H]\centering
\includegraphics[width=0.85\linewidth]{figures/exp3_traj_seed42_overview.png}
\caption{Seed 42 --- overview (median H vs.\ P loss trajectory, rank-aligned, IQR band).}
\end{figure}
\begin{figure}[H]\centering
\includegraphics[width=\linewidth]{figures/exp3_traj_seed42_grid.png}
\caption{Seed 42 --- per-equation trajectory grid (real, unaggregated H/P curves).}
\end{figure}
\begin{figure}[H]\centering
\includegraphics[width=0.68\linewidth]{figures/exp3_traj_seed42_winloss_scatter.png}
\caption{Seed 42 --- final-loss win/loss scatter.}
\end{figure}
\begin{figure}[H]\centering
\includegraphics[width=0.85\linewidth]{figures/exp3_traj_seed42_speedup_bar.png}
\caption{Seed 42 --- per-equation speedup to $\Rsq\ge0.9999$.}
\end{figure}

\FloatBarrier
\section{Cross-Seed Aggregate}
\label{sec:allseeds}

All 5 seeds $\times$ 12 equations pooled (60 equation--seed pairs per metric).

\begin{figure}[H]\centering
\includegraphics[width=0.75\linewidth]{figures/exp3_traj_ALLSEEDS_winloss_scatter.png}
\caption{Pooled final-loss win/loss across all 5 seeds, with paired Wilcoxon signed-rank test on
$\log_{10}(\text{P}) - \log_{10}(\text{H})$.}
\end{figure}

\begin{figure}[H]\centering
\includegraphics[width=0.92\linewidth]{figures/exp3_traj_ALLSEEDS_speedup_strip.png}
\caption{Pooled speedup to $\Rsq\ge0.9999$ across all 5 seeds (jittered points per equation, black
tick = per-equation median), with paired Wilcoxon signed-rank test on $\log(\text{speedup})$.}
\end{figure}

\begin{table}[H]
\centering
\caption{Cross-seed paired Wilcoxon signed-rank tests ($n=60$ equation--seed pairs each).}
\begin{tabular}{lrrrrr}
\toprule
Metric & Median diff. & Statistic & $p$-value & Favors H & Favors P \\
\midrule
$\log_{10}(\text{P}) - \log_{10}(\text{H})$ final loss & $-2.997$ & $153$ & $2.03\times10^{-8}$ & 8 & 52 \\
$\log(\text{P iters} / \text{H iters})$ to threshold & $1.195$ & $0$ & $1.63\times10^{-11}$ & 60 & 0 \\
\bottomrule
\end{tabular}
\end{table}

\FloatBarrier
\section{Conclusions}
\label{sec:conclusions}

\begin{enumerate}
  \item \textbf{The two candidate metrics disagree, and that disagreement is diagnostic, not
    noise.} Final loss says P wins in the large majority of equations; threshold-crossing speed
    says H wins in every single one. Both results are statistically strong
    ($p < 10^{-7}$ in both directions). This is expected given the synthetic data's construction
    (H's runs were made shorter but biased to cross the threshold earlier) and it illustrates a
    real risk for the actual paper: an unqualified ``H achieves lower loss'' or ``H is Nx faster''
    claim is only meaningful once the comparison protocol (matched iterations? matched wall-clock?
    time-to-threshold?) is stated explicitly.
  \item \textbf{Recommended primary metric: iterations (or wall-clock time) to reach the paper's
    own stated success threshold ($\Rsq \ge 0.9999$).} This is the metric least confounded by
    arms simply being allowed to run for different lengths, and it directly answers the practical
    question a warm-start is supposed to answer: ``does the LLM proposal let PySR reach the same
    target faster.''
  \item \textbf{Per-equation consistency matters as much as the pooled average.} The speedup strip
    plot (Fig.\ in Sec.~\ref{sec:allseeds}) shows the effect held up per-equation across all 5
    seeds in this run, not just as an average dominated by a few equations --- that per-equation
    view is what should be checked on the real data before citing a single pooled number.
  \item \textbf{This report is a pipeline validation, not a result.} Every figure and statistic
    here must be regenerated from the real \texttt{exp3\_nguyen12\_seed*.json} files (by pointing
    \texttt{--results-dir} at them) before any of these numbers are used in the paper. Given the
    separately-flagged Nguyen-12 success-count arithmetic issue (three disagreeing denominators
    for a nominally single success rate: 10/12, 11/12, 4/12), we recommend the same
    regenerate-and-cross-check discipline for every summary statistic in the paper, not only the
    ones covered by this pipeline: compute it once, programmatically, from source data, and have
    the generating script itself printed/exported the reconciled numbers used in the table or
    caption, rather than adjusting printed totals by hand.
\end{enumerate}

\end{document}
